\subsection{Loading-dependent pitch stability and a researcher-generated tail revision}
\label{sec:tail-loading-revision}

The small positive nominal margin is not representative of the entire declared
loading envelope. We independently re-summed all eight V15 loading cases,
including pilot masses of 65, 75 and 85 kg, the specified seat positions, and
full or empty fuel. The fuel tank structure remains present when its fuel is
removed. The mass ledger is still assumed rather than measured. Pilot intrinsic
inertia was scaled with mass at unchanged assumed body shape; changes in human
body proportions are not represented. Intrinsic products of inertia were not
provided for the components. Consequently, a sum of centroid parallel-axis
products must not be presented as a complete physical inertia tensor.

A first screen held the nominal neutral point fixed while moving the loading
CG. The two aft-CG cases were then independently re-trimmed, rather than accepting
that extrapolation as a stability result. Both use the 65 kg pilot at seat
coordinate $x=0.35$ m, with 8 kg or zero fuel. The original V15 fine-grid results
in Table~\ref{tab:tail-loading} have positive $C_{m_\alpha}$ and negative
$SM=-C_{m_\alpha}/C_{L_\alpha}$, although their lift and pitch moment satisfy
the surface-model balance. This provides a direct example of why equilibrium
and stability are distinct: the controls can establish a stationary operating
point while a small angle-of-attack disturbance produces a destabilizing
pitching-moment increment. Both mesh levels predict the same sign. The result
does not establish the behavior of the installed aircraft, because the body,
propulsion and other omitted effects remain unresolved.

\begin{table}[htbp]
\centering\small
\caption{Fine-grid loading comparison at 13 m/s. The revised geometry is researcher-generated, not a new Fable response. All margins are surface-model predictions.}
\label{tab:tail-loading}
\begin{tabular}{llrrrr}
\toprule
Geometry & Fuel (kg) & Mass (kg) & $\alpha$ (deg) & $C_{m_\alpha}$ (rad$^{-1}$) & $SM$ (\%)\\
\midrule
Original V15 & 8 & 343.2 & 5.31818 & +0.057983 & -1.6703\\
Original V15 & 0 & 335.2 & 5.02707 & +0.055162 & -1.5849\\
Enlarged tail & 8 & 349.2 & 5.25153 & -0.14426 & +3.9463\\
Enlarged tail & 0 & 341.2 & 4.96037 & -0.14354 & +3.9166\\
Tail plus test mass & 0 & 343.2 & 4.94835 & -0.09612 & 2.62295\\
\bottomrule
\end{tabular}
\end{table}

We next made an explicitly researcher-generated design experiment. This step
must not be attributed to the language model or counted as a successful new
model response. The root station, chord and hinge location were retained,
while the exposed span length of each horizontal-tail half was multiplied by
$k$. With the original root at $|y|=0.55$ m and half-panel length 2.30 m, the
new outer station is $|y|=0.55+2.30k$. Thus $S_t=4.60k$ m$^2$; the overall
span including the center gap is $1.10+4.60k$ m, not $5.70k$. A constant
assumed areal mass gives $m_t=8k$ kg. Each half-panel centroid and thin-plate
intrinsic inertia were recomputed before recalculating aircraft mass and CG.
The added tail mass moves the CG aft and therefore partly offsets the
aerodynamic gain. No additional reinforcement or actuator mass is claimed to
be zero: those unresolved contributions are excluded from this conditional
trade and tested separately below.

\begin{table}[htbp]
\centering\small
\caption{Tail-span trade for the light-pilot, full-fuel case on the coarse mesh.}
\label{tab:tail-span-trade}
\begin{tabular}{rrrrr}
\toprule
$k$ & Tail area (m$^2$) & Aircraft mass (kg) & Absolute tail angle (deg) & $SM$ (\%)\\
\midrule
1.15 & 5.29 & 344.4 & -4.1456 & -0.4558\\
1.30 & 5.98 & 345.6 & -2.5306 & +0.6693\\
1.50 & 6.90 & 347.2 & -1.0228 & +2.1721\\
1.75 & 8.05 & 349.2 & +0.2244 & +4.1450\\
\bottomrule
\end{tabular}
\end{table}

An exploratory 3\% margin was used to select a candidate for further study;
it is an evaluator choice, not an airworthiness standard, handling-quality
requirement or validated uncertainty allowance. The first three candidates
did not meet this screen. The $k=1.75$ candidate has area 8.05 m$^2$, overall
span 9.15 m and assumed tail mass 14 kg. Its nominal all-up mass is 359.2 kg.
All eight coarse-grid loading states were re-trimmed, yielding positive
margins from 4.115\% to 9.713\%. The two critical fine-grid margins are
3.946\% and 3.917\%; their absolute tail angles are $+0.28375^\circ$ and
$+0.32297^\circ$. The coarse--fine difference is approximately 0.20 percentage
points in these cases. This is a limited mesh comparison, not a complete
convergence demonstration. The remaining six fine-grid loading cases and
model-form uncertainty still require evaluation. A positive margin alone
does not establish acceptable phugoid, short-period, Dutch-roll, roll or spiral
behavior, nor does it prove usable control forces or recovery capability.

\paragraph{Structural mass feedback rather than a free aerodynamic improvement.}
Enlarging the tail increases its structural demands. For a simple cantilever
of length $l$ under unchanged uniform line load $w$, the root moment is
$wl^2/2$ and the tip deflection is $wl^4/(8EI)$, where $E$ is Young's modulus
and $I$ is the second moment of area. Holding $w$ and $EI$ fixed gives scaling
factors $k^2=3.0625$ and $k^4=9.3789$ at $k=1.75$. These are illustrative
beam scalings, not computed loads or deflections for this aircraft: the actual
span loading, spar section, joint stiffness and maneuver conditions have not
been established. Root fittings, bearings, torque transmission, control
actuation, swept clearance and failure loads must be sized before accepting
the revision. Its drag and mass changes also invalidate any automatic reuse
of the original power margin.

To quantify one consequence of that open structural loop, an explicit
2 kg point mass was added at the central tail hinge, $(5.65,0,0.775)$ m,
and the light-pilot, empty-fuel case was re-trimmed on the fine mesh. This is
a sensitivity experiment, not an estimate of sufficient reinforcement.
Zero intrinsic inertia is used only for this mathematical point mass; real
hardware inertia is unknown. The resulting margin is 2.62295\%, compared
with 3.9166\% before the added mass. A complementary fixed-neutral-point
screen estimates the mass increment that consumes a chosen margin:
\begin{equation}
 \Delta m = \frac{M(x_{\mathrm{lim}}-x_{\mathrm{CG}})}
 {x_r-x_{\mathrm{lim}}},\qquad
 x_{\mathrm{lim}}=x_{\mathrm{NP}}-\bar c\,SM_{\mathrm{target}}.
\label{eq:tail-mass-budget}
\end{equation}
Here $M$ and $x_{\mathrm{CG}}$ are the pre-increment mass and CG, $x_r$ is the
added-mass station, $x_{\mathrm{NP}}$ is the frozen model neutral point and
$\bar c=1.85$ m is the reference chord. For the empty-fuel case the screen
gives approximately 1.45000 kg before reaching the exploratory 3\% margin.
Because re-trim changes aerodynamic derivatives, this is a screening budget,
not an allowable installed mass. The next design iteration should compare
structurally sized tail variants and feasible forward mass redistribution,
then repeat geometry, mass, power, trim and derivative checks. The current
candidate closes neither the structural nor the dynamic-mode acceptance gate.

All original model responses and superseded evaluator results are retained.
The loading audit, tail trade and reinforcement sensitivity have separate
scripts, raw solver inputs/outputs and checksums in the local review package.
They are not yet part of the immutable public release identified in the
data-availability statement.

